\section{Measurement Methods}
\label{sec:methods}

Throughout, $x\in\mathbb{R}^{T}$ is a context, $y\in\mathbb{R}^{H}$ its continuation, $\hat y=f(x)$ a point forecast and
$h^{(\ell)}$ the residual stream after block $\ell$. The analyses are organized in levels: L0 (behavior), L1
(geometry), L2 (linear translatability) and L3 (perturbation and patching), followed by depth analysis and causal
concepts.

\subsection{Behavior, geometry and perturbation (L0--L3)}
\label{sec:l0}
\label{sec:l1l2}
\label{sec:l3}

\paragraph{L0: behavior.} Accuracy is the per-series MASE \citep{hyndman2006mase}, with the seasonal lag taken from
ground truth and a floor on the naive denominator; without the floor, intermittent series spuriously improve
(\fid{MN-23}). sMAPE, pinball loss \citep{gneiting2007scoring} and interval calibration are reported alongside. Model
pairs are compared per generator family with a paired series bootstrap, Holm-corrected across all pairs and families
\citep{holm1979}, and differences are read against each model's repeat-run noise floor. Two further stages measure
parameters, FLOPs, latency and the share of a forecast's FLOPs spent in captured blocks, and test what each model's
input pipeline does before any layer runs (scale equivariance, handling of missing values, quantization).

\paragraph{L1 and L2: geometry.} For aligned, window-pooled and centered activations $A$ and $B$ over the same series,
linear CKA \citep{kornblith2019cka},
$\mathrm{CKA}(A,B)=\|B^{\top}A\|_F^2/(\|A^{\top}A\|_F\|B^{\top}B\|_F)$, is computed for every layer pair of every
model pair with a cluster bootstrap and a shuffled-series null. L2 fits ridge maps between layers
\citep{bansal2021stitching} and reports only the \emph{gain} over a baseline that predicts the target layer from input
features, since much of any layer is linearly predictable from the input alone. Agreement between activation clusterings
is measured by adjusted mutual information \citep{vinh2010ami}.

\paragraph{L3: perturbation.} For each corruption $c$, the sensitivity \emph{fingerprint} of a model is the relative
activation change $\|h^{(\ell)}(c(x))-h^{(\ell)}(x)\|/\|h^{(\ell)}(x)\|$ as a function of layer, and its
\emph{restoration curve} is the fraction of forecast damage undone by patching clean activations of layer $\ell$, whole
or per window, into the corrupted pass. Two models are compared by the rank correlation of their fingerprints on a
shared relative-depth axis, never by transplanting activations between models. Restoration is causal within a model;
agreement between two models' fingerprints is correlational.

\subsection{Depth}
\label{sec:depth}

A skip lens decodes the forecast from each layer through the model's own head, and a tuned ridge lens
\citep{belrose2023tuned} fits a readout per layer. The depth axis is the fraction of the \emph{whole} stack, including
parts that are not captured: Chronos-T5's decoder is not captured, so its last captured encoder block sits at depth
0.478 (\fid{DE-04}). The lens is withheld for a model whose head reads positions that patching never writes, where it
would return a flat and meaningless curve (\fid{DE-03}).

\subsection{Causal concepts}
\label{sec:concepts}

\paragraph{Sparse dictionaries.} Each target layer receives a TopK sparse autoencoder (SAE)
\citep{makhzani2013ksparse,gao2025topk},
\begin{equation}
  z(h)=\operatorname{TopK}\big(W_e(h-b_d)+b_e\big), \qquad \hat h=W_d\, z(h)+b_d ,
\end{equation}
with a searched dictionary size, an auxiliary loss that revives dead atoms and a minimum number of training steps
(\fid{MP-03}). Because single SAE seeds are not reproducible \citep{paulo2025seeds}, three replicate dictionaries per
target provide a within-model reproducibility ceiling (\fid{MP-04}).

\paragraph{The ablation battery.} Feature $i$ is tested on the $k$ series on which it fires most strongly. The
forecast under the reconstruction with $z_i=0$, written $\hat h_{-i}$, is compared with the forecast under the full
reconstruction:
\begin{equation}
  \Delta_i^{(c)}(x) = g_c\big(f(x;h^{(\ell)}\!\leftarrow\!\hat h_{-i})\big) - g_c\big(f(x;h^{(\ell)}\!\leftarrow\!\hat h)\big),
  \label{eq:ablation}
\end{equation}
where $g_c$ is one of nine forecast \emph{channels}: trend, seasonal strength, spectral centroid, level, dispersion,
near- and far-horizon shape distance, MASE and flatness. Using the reconstruction as the baseline keeps the SAE's
reconstruction error in both arms. This matters because that error moves the forecast several times more than a typical
feature does (Section~\ref{sec:rq3}). The two shape channels are distances and are therefore treated as unsigned.

\paragraph{Random-direction nulls.} The null asks whether removing the feature's direction matters more than removing
the same amount of activation along another direction. Two variants are used.
\begin{itemize}
  \item The \emph{profile-matched} null (used for the v2 corpus) removes, on the same rows and at every token $t$,
    exactly the amount the feature removes, $|z_i(t)|\,\|w_i\|$, along each of 16 random \emph{decoded} directions: a
    unit vector drawn isotropically in the dictionary's code space and mapped through the decoder. Only the direction
    differs from the real ablation. We call this null isotropic for short.
  \item The \emph{covariance-matched} null (used for the v3 corpus) is identical except for the direction, which is a
    random Gaussian combination of the same rows' centered clean token activations. It is thus a draw from a Gaussian
    with the layer's sample activation covariance and lies in the subspace the model actually uses. It is the stricter
    of the two: on an earlier four-model run it reduced the clearing cells from 3.06 times the empirical chance rate to
    1.13 times (Appendix~\ref{app:faq}).
\end{itemize}
Feature $i$ \emph{clears} channel $c$ if its mean $|\Delta_i^{(c)}|$ exceeds the 95th percentile $q_{95}^{(c)}$ of the
pooled null, and its effect vector is $v_i^{(c)}=\overline{\Delta_i^{(c)}}/q_{95}^{(c)}$, the signed effect in units of
its own null. Because a mean over rows is compared with a percentile of single rows, the chance rate of clearing is not
5\%. It is estimated per cell by scoring each null draw against the others (leave-one-draw-out), and ranged from 0.3\% to
4\% per target in the case study. Matching the removed amount token by token is what lets the null pass the known-answer
test of Section~\ref{sec:knownanswer}; a simpler null that removes one uniform, batch-averaged amount fails it
(\fid{MN-29}).

\paragraph{Reach controls.} Before any feature is scored, a \emph{reach probe} requires that patching a layer's own
activations back into it changes the forecast by exactly $0.0$, and that patching another layer's activations changes it
by a relative amount above $10^{-3}$. A target that fails either check is withheld with its reason
(Figure~\ref{fig:battery}).

\begin{figure}[tb]
\centering
\begin{tikzpicture}[
  box/.style={draw, rounded corners=2pt, align=center, font=\scriptsize, minimum height=9mm, fill=white, text width=21mm},
  gate/.style={draw, diamond, aspect=2.0, align=center, font=\scriptsize, inner sep=0.5pt, fill=soft, text width=15mm},
  arr/.style={-{Stealth[length=1.6mm]}, gray!70!black},
  lab/.style={font=\scriptsize\itshape, text=gray!70!black}]
  \node[box] (tgt) at (0,0) {target layer $\ell$\\+ trained SAE};
  \node[gate] (self) at (3.5,0) {self-patch\\$=0.0$?};
  \node[gate] (cross) at (7.0,0) {cross-patch rel.\\$>10^{-3}$?};
  \node[box] (rows) at (10.6,0) {feature $i$: its\\top-$k$ firing series};
  \node[box, fill=bad!8] (wh) at (5.25,1.7) {target withheld,\\reason rendered};
  \node[box] (abl) at (10.6,-2.2) {forecast with $\hat h_{-i}$\\vs.\ with $\hat h$\\(same seed)};
  \node[box] (null) at (7.0,-2.2) {row-matched null:\\random direction, the\\feature's own profile};
  \node[box] (chan) at (3.5,-2.2) {each of 9 channels:\\$|\overline{\Delta}|>q_{95}$ of null?};
  \node[box, fill=accent!12] (vec) at (0,-2.2) {causal feature\\$v_i=\overline{\Delta}/q_{95}$};
  \draw[arr] (tgt) -- (self);
  \draw[arr] (self) -- node[above,lab]{yes} (cross);
  \draw[arr] (cross) -- node[above,lab]{yes} (rows);
  \draw[arr] (self.north) |- node[left,lab,pos=0.3]{no} (wh.west);
  \draw[arr] (cross.north) |- node[right,lab,pos=0.3]{no} (wh.east);
  \draw[arr] (rows) -- (abl);
  \draw[arr] (abl) -- (null);
  \draw[arr] (null) -- (chan);
  \draw[arr] (chan) -- node[below,lab]{$\ge$1 clears} (vec);
\end{tikzpicture}
\caption{The ablation battery. A target is scored only after both reach controls pass; each feature is compared,
channel by channel, with a random-direction null on exactly its own top-firing series.}
\label{fig:battery}
\end{figure}

\paragraph{From features to concepts.} Causal features are grouped at three granularities.
\begin{itemize}
  \item \emph{Per-target concepts} cluster one layer's effect vectors with $k$-means, choosing $k$ by silhouette
    \citep{rousseeuw1987silhouette} among values whose smallest cluster has at least three members. Without this
    constraint the best silhouettes are one-outlier splits (\fid{SH-20}). A target with no admissible $k$ is called
    \emph{non-modular}.
  \item The \emph{atlas} pools features from all models and layers and cuts a complete-linkage dendrogram at cosine
    0.9, so an atlas concept is a tight, shared causal-effect profile.
  \item \emph{Families} are a coarser tiling for reading (average linkage at cosine 0.5), tested against a
    column-shuffle null and titled deterministically from their channels.
\end{itemize}

\subsection{Two notions of sharing}
\label{sec:sharing}

Whether a concept is ``shared'' between two models is two separate questions.

\paragraph{Same inputs.} \emph{Concept transfer} asks whether some feature in another model separates a concept's
top-$k$ series from the rest, measured by AUC, against stratum-matched random subsets of series. It is run in both
directions; a transfer is \emph{reciprocal} if both directions clear, with Benjamini--Hochberg control
\citep{benjamini1995fdr} per ordered model pair and direction. Input agreement between two concept parts is measured by
the overlap of their top-$k$ series (hypergeometric test), never by whole-series correlation: two sparse features can
correlate at $\rho=0.80$ while sharing no top series (\fid{MN-19}).

\paragraph{Same effect on the same inputs.} For every reciprocal transfer, both sides are ablated on the same series set
$U$. Each side must first clear its own null on $U$. The two sides then agree if their level concordance and shape
cosine both exceed the floors given by ablating random feature sets matched to each side's activation. A statistic
below both floors' 5th percentile is called \emph{acts differently}; one that merely fails to exceed the 95th
percentile is \emph{no specific agreement}. The distinction matters: reading the second case as the first once labeled
109 of 288 tests as disagreements, against 9 under the correct test (\fid{MN-18}).

\paragraph{Evidence ladder and sharing classes.} The report summarizes each atlas concept on a six-rung ladder: R1 same
forecast effect (atlas membership), R2 same inputs, R3 reproducible across SAE seeds, R4 other dictionaries select the
same inputs, R5 same causal effect on the same inputs, and R6 confirmed on private data (the report labels the rungs L1--L6). Each concept also receives a
\emph{sharing class}: \emph{shared} (same effect and same inputs), \emph{partially shared}, \emph{convergent} (same
effect reached on different inputs) or \emph{single-model}. A concept part is flagged as provenance-driven when at least
80\% of its top series come from one real-derived generator.
